\par\medskip
\begin{center}\hypertarget{native-164-article117}{}\pdfbookmark[2]{117.}{article-117}117.\end{center}
Der auf diese Weise aus den Tafeln berechnete Ort der Erde ist auf den Erdmittelpunkt zu beziehen, der beobachtete Ort des Himmelskörpers hingegen auf einen Punkt der Erdoberfläche. Diese Abweichung lässt sich auf drei Arten beheben. Man kann die Beobachtung auf den Erdmittelpunkt reduzieren, sie also von der Parallaxe befreien. Man kann den heliozentrischen Ort der Erde auf den tatsächlichen Beobachtungsort reduzieren, indem man die Parallaxe auf den aus den Tafeln berechneten Ort der Sonne richtig anwendet. Schließlich kann man beide Orte auf einen dritten Punkt beziehen, den man am bequemsten als Schnittpunkt der Sichtlinie mit der Ekliptikebene wählt. Die Beobachtung selbst bleibt dabei unverändert. Wie man den Ort der Erde auf diesen Punkt reduziert, haben wir in Artikel 72 gezeigt. Das erste Verfahren lässt sich nur anwenden, wenn der Abstand des Himmelskörpers von der Erde wenigstens annähernd bekannt ist. Dann ist es recht bequem, besonders wenn die Beobachtung unmittelbar im Meridian stattfindet, wo sich die Parallaxe nur auf die Deklination auswirkt. Im Übrigen ist es vorzuziehen, diese Reduktion unmittelbar am beobachteten Ort vorzunehmen, bevor man die Umrechnungen des vorhergehenden Artikels ausführt. Ist der Abstand von der Erde jedoch noch völlig unbekannt, so muss man auf das zweite oder dritte Verfahren zurückgreifen. Das zweite kommt zur Anwendung, wenn der Äquator als Grundebene gewählt wird. Das dritte ist vorzuziehen, wenn man alle Orte auf die Ekliptik beziehen will.
